% SPDX-License-Identifier: CC-BY-SA-4.0
% Author: Matthieu Perrin

\begingroup

\begin{exercice}[Expressions rationnelles correspondant à un automate]\label{exo:lexing/kleene/automaton_to_regexp}

  Traduire chaque automate ci-dessous en une expression rationnelle sur l'alphabet $\{a,b\}$.
  \begin{question}
  \item
    \begin{tikzpicture}[automaton, size=20mm, baseline=(1.base)]
      \state[initial]   (1) at (0,0) {$1$}; 
      \state            (2) at (1,0) {$2$}; 
      \state[accepting] (3) at (2,0) {$3$};
      
      \path (1) edge             node {$a$}    (2);
      \path (2) edge             node {$b$}    (3);
      \path (1) edge[loop above] node {$a, b$} (1);
      \path (3) edge[loop above] node {$a, b$} (3);
    \end{tikzpicture}

  \item
    \begin{tikzpicture}[automaton, size=20mm, baseline=(1.base)]
      \state[initial, accepting]    (1) at (0,0) {$1$}; 
      \state                        (2) at (1,0) {$2$};
      
      \path (1) edge[bend left]  node {$a$} (2);
      \path (2) edge[bend left]  node {$b$} (1);
      \path (1) edge[loop above] node {$b$} (1);
    \end{tikzpicture}

  \item
    \begin{tikzpicture}[automaton, size=20mm, baseline=(1.base)]
      \state[initial, accepting] (1) at (0,1) {$1$}; 
      \state                     (2) at (1,1) {$2$}; 
      \state                     (3) at (1,0) {$3$}; 
      \state[accepting]          (4) at (2,1) {$4$}; 

      \path (1) edge[loop above] node       {$b$} (1);
      \path (2) edge[loop above] node       {$a$} (2);
      \path (1) edge             node       {$a$} (2);
      \path (2) edge             node[swap] {$b$} (3);
      \path (3) edge             node       {$b$} (1);
      \path (4) edge             node[swap] {$a$} (2);
      \path (3) edge             node       {$a$} (4);
      \path (4) edge[bend left]  node       {$b$} (3);
    \end{tikzpicture}    
  \end{question}
  
\end{exercice}

\begin{correction}

  Il faut passer par le système d'équations rationnelles, puis utiliser le lemme d'Arden.
  
  \begin{question}
  \item $\left\{\begin{array}{lllll}
    L_1 & = & a L_1 \mid b L_1 \mid a L_2  & = & (a\mid b)L_1 \mid a L_2\\
    L_2 & = & b L_3 \\
    L_3 & = & a L_3 \mid b L_3 \mid \varepsilon & = & (a\mid b)L_3 \mid \varepsilon\\
  \end{array}\right.$
    
    \begin{itemize}
    \item $L_3 = (a\mid b)L_3 \mid \varepsilon$. $\varepsilon\notin\{a, b\}$ donc d'après le lemme d'Arden, $L_3 = (a\mid b)^\star$ 
    \item $L_2 = b L_3 = b (a\mid b)^\star$
    \item $L_1 = (a\mid b)L_1 \mid a L_2 = (a\mid b)L_1 \mid a b (a\mid b)^\star$. $\varepsilon\notin\{a, b\}$ donc d'après le lemme d'Arden, $L_1 = (a\mid b)^\star ab (a\mid b)^\star$.
    \end{itemize}
    
  \item $\left\{\begin{array}{lll}
    L_1 & = & a L_2 \mid b L_1 \mid \varepsilon\\
    L_2 & = & b L_1 \\
  \end{array}\right.$
    
    \begin{itemize}
    \item $L_1 = a L_2 \mid b L_1 \mid \varepsilon = a b L_1 \mid b L_1 \mid \varepsilon = (a b \mid b) L_1 \mid \varepsilon$.
      $\varepsilon\notin\{ab, b\}$ donc d'après le lemme d'Arden, $L_1 = (ab \mid b)^\star$.
    \item $L_2 = b L_1 = b(ab \mid b)^\star$
    \end{itemize}

  \item
    $\left\{\begin{array}{lllll}
    L_1 & = & a L_2 \mid b L_1 \mid\varepsilon &=& (b^\star a^+ b)^+ (a \mid b^+) \mid b^\star \\ 
    L_2 & = & a L_2 \mid b L_3                 &=& a^\star b (b^\star a^+ b)^\star (a \mid b^+)\\ 
    L_3 & = & a L_4 \mid b L_1                 &=& (b^\star a^+ b)^\star (a \mid b^+) \\ 
    L_4 & = & a L_2 \mid b L_3 \mid\varepsilon &=& a^\star b (b^\star a^+ b)^\star (a \mid b^+) \mid \varepsilon\\ 
  \end{array}\right.$
    
    \begin{itemize}
    \item $L_2 = a L_2 \mid b L_3$.  $\varepsilon\notin\{a\}$ donc d'après le lemme d'Arden, $L_2 = a^\star b L_3$.
    \item $L_1 = a L_2 \mid b L_1 = b L_1 \mid a^+ b L_3 \mid\varepsilon$. On a $\varepsilon\notin\{b\}$ donc d'après le lemme d'Arden, $L_1 = b^\star a^+ b L_3 \mid b^\star$.
    \item $L_4 = a L_2 \mid b L_3 \mid\varepsilon = a^+ b L_3 \mid b L_3 \mid\varepsilon = a^\star b L_3 \mid \varepsilon$
    \item $L_3 = a L_4 \mid b L_1 = a^+ b L_3 \mid a \mid b^+ a^+ b L_3 \mid b^+ = b^\star a^+ b L_3 \mid a \mid b^+  = (b^\star a^+ b)^\star (a \mid b^+)$
    \end{itemize}

  \end{question}

\end{correction}

\endgroup
\endinput
